\Needspace{8\baselineskip}\originalchapter{2003 作业四: 碰撞、自由程与逃逸}
题源: 2 Oct 2003, due 9 Oct. Coulomb对数依截止约定有O(1)差, 以下给可复算的一套约定.
\begin{exercise}\label{ex:03-4-1}
原题1. ne=ni=$10^{20}$m$^{-3}$、Te=5keV. 求ln\ensuremath{\Lambda}: (a)100keV电子束撞静止氢离子; (b)同束撞等离子体电子; (c)热碰撞.
\end{exercise}
\begin{solution}
取最大冲量距离\ensuremath{\lambda}De=$5.256\times10^{-5}$m; 最小距离为 $b_{\min}=\max[e^2/(4\pi\epsilon_0\mu v^2),\hbar/(\mu v)]$, $\ln\Lambda=\ln(\lambda_D/b_{\min})$. 首先按题的非相对论 $v=\sqrt{2E/m_e}$ 与静止靶相对速, (a)\ensuremath{\mu}\ensuremath{\simeq}me, 得18.26; (b)\ensuremath{\mu}=me/2, 得17.57, 忽略Te/E小修正; (c)取E=Te、\ensuremath{\mu}=me/2为e-e代表, 得16.07. 若取e-i热碰撞则约16.76. 100keV已约0.2me c\ensuremath{{}^{2}}, 精确v与量子动量使数值有小改动; 这里不把以上约定说成唯一精确碰撞对数. 直接在100keV用经典b90而不检查量子截止会高估.
\end{solution}
\begin{exercise}\label{ex:03-4-2}
原题2. 分布 $f=n[\pi^2v_t^3]^{-1}[1+(\mathbf u-\mathbf u_d)^2]^{-2}$, $\mathbf u=\mathbf v/v_t$, $\mathbf u_d=u_d\mathbf e_x$, ud\ensuremath{\ll}1. 一阶展开, 代撞静止靶的动量损失, 用各向同性化速度积分, 求\ensuremath{\nu}ei.
\end{exercise}
\begin{solution}
展开 $f=f_0+4nu_du_x/[\pi^2v_t^3(1+u^2)^3]$. 快电子重离子阻力 $\dot{\mathbf v}=-K\mathbf v/v^3$, $K=n_ie^4\ln\Lambda/(4\pi\epsilon_0^2m_e^2)$. 只有奇的一阶项贡献x动量. 角积分 $\int u_x^2\dd\Omega=4\pi u^2/3$, 径向 $\int_0^\infty u(1+u^2)^{-3}\dd u=1/4$, 得
\begin{equation}
 \frac1n\int\frac{v_x}{v^3}f\dd^3v=\frac{4u_d}{3\pi v_t^2},\quad
 \nu_{ei}=\frac{4K}{3\pi v_t^3}=\frac{n_ie^4\ln\Lambda}{3\pi^2\epsilon_0^2m_e^2v_t^3}.
\end{equation}
此v$^{-4}$三维尾的绝对一阶矩与能量不有限; ``漂移''可按分布中心、对称主值或物理高速度截止理解. 阻力积分收敛, 但不把它变为有有限热能的Maxwell等效温度. 静止重靶近似也不能用于任意电子相对速.
\end{solution}
\begin{exercise}\label{ex:03-4-3}
原题3. 自由程定义为特征速度除动量损失频率. 求 (a)氢tokamak, Te=Ti=10keV,n=$.5\times10^{20}$m$^{-3}$的热电子; (b)同热离子; (c)加工等离子体Te=5eV,n=$5\times10^{18}$m$^{-3}$热电子.
\end{exercise}
\begin{solution}
明确取 $v_t=\sqrt{2T/m}$, 常用热频率近似 $\nu_e=2.91\times10^{-6}n_{\rm cm^{-3}}\ln\Lambda/T_{\rm eV}^{3/2}$, $\nu_i=4.80\times10^{-8}n_{\rm cm^{-3}}\ln\Lambda/T_{\rm eV}^{3/2}$ (氢). \ensuremath{\lambda}De/最大经典量子截止给tokamak ln\ensuremath{\Lambda}\ensuremath{\simeq}17.80、加工\ensuremath{\simeq}10.85. 得(a) $2.29\times10^4$m、(b) $3.24\times10^4$m、(c).0939m. 频率定义与热速取法会改变O(1)系数, 自由程长于装置说明局部无碰撞, 不是粒子完全无散射或无约束.
\end{solution}
\begin{exercise}\label{ex:03-4-4}
原题4. 圆截面环形氢等离子体Te=2keV均匀,a=.3m,R=1.2m, 驱动I=$4\times10^5$A. 最低a/R阶、可取generic ln\ensuremath{\Lambda}, 求E\ensuremath{\phi}和单匝loop voltage. n=$10^{19}$m$^{-3}$, 忽略相对论求电子成为runaway的最小平行能, 再判断该忽略是否合理.
\end{exercise}
\begin{solution}
取ln\ensuremath{\Lambda}=17, $\eta\simeq5.2\times10^{-5}\ln\Lambda/T_{\rm eV}^{3/2}=9.883\times10^{-9}$\ensuremath{\Omega}m. $J=I/(\pi a^2)=1.415\times10^6$A/m\ensuremath{{}^{2}}, $E=\eta J=.01398$V/m, $V_{loop}=2\pi RE=.1054$V. 快电子同时撞离子及电子的非相对论减速系数采用 \begin{equation}F\simeq\frac{3ne^4\ln\Lambda}{4\pi\epsilon_0^2m_ev^2}.\end{equation} 令eE=F, 得 $W_{\min}=3ne^3\ln\Lambda/(8\pi\epsilon_0^2E)\simeq.4752$MeV. 形式上的非相对论速度$\sqrt{2W/m_e}\simeq1.364c$已超过光速. 该能量接近me c\ensuremath{{}^{2}}=.511MeV, 因而非相对论仅给模型估计, 不足以精确预测逃逸阈值. 只含离子阻力会少一部分, 所用碰撞模型必须说明.
\end{solution}
